\section{Turbomachinery Design}

The design point for this thesis was chosen from the capabilities of the test rig. The design parameters for the pump and turbine is shown in \todo{addcreftable} and this section explains how each number was chosen.

The RPM limit was set by the maximum speed of the Magtrol TM306 inline torquemeter and the AHB-3 hysteris brake at 20,000 RPM. The pressure rise and flow rate was chosen to study the low specific regime. The supply pump was decided to be Hozelock's JET 3000 K7 Jet pump as it was available in the department and had a published performance curve. At \valPumpMdotDesign, the supply pump had a head of around \valSupplyPumpHeadDesign. Then, the head rise of the pump was chosen to be \valPumpPressureDesign for a specific speed of \todo{addspecificspeed}.



% \todo{does lock's model change barske's design methodology for d2? I can write the barske model for now and say future work can look at rearranging lock maybe?}

\section{Turbine Sizing and Mechanical Design}


\subsection{Turbine Design Point}

Turbine sizing has to satisfy many requirements simultaneously, many design variables. After the pump was sized for \valPumpRpmDesign at a power of \valTurbPowerDesign, there are many remaining free variables.

The main parameters being turbine diameter $D_mean$, mass flow through the turbine $\dot{m}$, blade inlet angle $\beta_3$.

